\section{Recursive feasibility}
\label{sec:feasibility}

A controller that applies the first step of a plan needs a plan at every step. If the
optimization problem has no admissible solution at some time $k$, the controller has nothing
to apply. The property that rules this out is the subject of this chapter.

\subsection{Feasible sets and invariance}

\begin{definition}[feasible set]
The \emph{feasible set} $\mathcal{F}_N$ of problem \eqref{eq:PN} is the set of states $x$ from
which the problem admits at least one admissible input sequence.
\end{definition}

\begin{definition}[recursive feasibility]
The MPC algorithm is \emph{recursively feasible} if, whenever $\mathcal{P}_N(x_k)$ has a solution,
$\mathcal{P}_N(x_{k+1})$ has a solution as well, with $x_{k+1} = f(x_k, \kappa_N(x_k))$. In
words: feasible once, feasible forever.
\end{definition}

Recursive feasibility is a statement about sets. For an autonomous system $x_{k+1} = g(x_k)$ a
set $S$ is \emph{positively invariant} when every state in $S$ is mapped into $S$: $x \in S
\Rightarrow g(x) \in S$. For a controlled system $x_{k+1} = f(x_k,u_k)$ a set $S$ is
\emph{controlled invariant} when for every $x \in S$ there exists an admissible input $u \in \U$
such that $f(x,u) \in S$. Recursive feasibility is exactly the statement that the MPC law makes
the feasible set controlled invariant: from every state in $\mathcal{F}_N$, the input it chooses
leads to a state that is again in $\mathcal{F}_N$ (\cref{fig:invariant}).

\begin{figure}[t]
\centering
\begin{tikzpicture}[x=1cm,y=1cm,>=Stealth]
  \draw[thick,fill=mpc!10] plot[smooth cycle] coordinates {(0,0) (2.2,-0.6) (4.2,0.3) (4.6,2.2) (3.2,3.4) (1.0,3.2) (-0.6,1.8)};
  \draw[thick,fill=sun!25] plot[smooth cycle] coordinates {(1.6,1.2) (2.8,1.0) (3.3,1.9) (2.6,2.6) (1.6,2.3)};
  \node[mpc!70!black] at (0.6,0.5) {$\mathcal{F}_N$};
  \node[sun!70!black] at (2.45,1.8) {$\Xf$};
  % trajectory staying inside
  \fill (3.9,0.9) circle (0.06);
  \draw[->,thick] (3.9,0.9) -- (3.2,1.4);
  \draw[->,thick] (3.2,1.4) -- (2.7,1.75);
  \draw[->,thick] (2.7,1.75) -- (2.45,2.0);
  \node[right,font=\small] at (3.95,0.85) {$x_k$};
  \node[font=\small,align=left,anchor=west] at (5.3,2.6) {every arrow stays inside:\\the set is controlled invariant};
\end{tikzpicture}
\caption{The feasible set $\mathcal{F}_N$ and the terminal set $\Xf$. Recursive feasibility
means that the control law never takes the state out of $\mathcal{F}_N$; the terminal set is
the smaller set, controlled invariant under the terminal law, that makes the proof work.}
\label{fig:invariant}
\end{figure}

\subsection{The proof}

The proof is constructive, and the construction is the natural one. Assume the nominal
conditions, so that the state at $k+1$ is exactly the one the controller predicted,
$x_{k+1} = x^*_{\ik{1}}$. At time $k$ the controller found the optimal sequence
$\{u^*_{\ik{0}}, u^*_{\ik{1}}, \dots, u^*_{\ik{N-1}}\}$ with predicted states
$\{x^*_{\ik{0}}, \dots, x^*_{\ik{N}}\}$, and $x^*_{\ik{N}} \in \Xf$. At time $k+1$ build the
\emph{tentative} sequence
\begin{equation}
  \tilde\bu_{k+1} = \{\,u^*_{\ik{1}},\, u^*_{\ik{2}},\, \dots,\, u^*_{\ik{N-1}},\,
  \kappa_f(x^*_{\ik{N}})\,\},
  \label{eq:tail}
\end{equation}
the tail of the previous optimal sequence, completed with one more input given by a
\emph{terminal control law} $\kappa_f$ applied to the previous terminal state. The first $N-1$
inputs are admissible and produce admissible states, because they did so at time $k$ --- the
window has merely moved. Two things remain to be checked: that the appended input is
admissible, $\kappa_f(x^*_{\ik{N}}) \in \U$, and that the state it produces satisfies the
terminal constraint, $f\big(x^*_{\ik{N}}, \kappa_f(x^*_{\ik{N}})\big) \in \Xf$.

Both hold at once if the terminal set and the terminal law are chosen so that
\begin{equation}
  x \in \Xf \ \Rightarrow\ \kappa_f(x) \in \U \ \text{ and }\ f\big(x,\kappa_f(x)\big) \in \Xf ,
  \label{eq:terminalinv}
\end{equation}
that is, if $\kappa_f$ makes $\Xf$ positively invariant for the controlled system while
respecting the input constraints. Then \eqref{eq:tail} is an admissible sequence for
$\mathcal{P}_N(x_{k+1})$, the problem has at least one solution, and recursive feasibility
follows.

\begin{theorem}[recursive feasibility]
Under the nominal assumptions, if $\Xf$ and $\kappa_f$ satisfy \eqref{eq:terminalinv}, then
$\mathcal{P}_N(x_k)$ feasible implies $\mathcal{P}_N(x_{k+1})$ feasible, and the feasible set
$\mathcal{F}_N$ is controlled invariant under the MPC law.
\end{theorem}

\subsection{Choosing the terminal set}

Two examples of terminal law appeared in class. The law that does nothing, $\kappa_f(x) = 0$,
works when $\Xf$ is invariant for the uncontrolled system. The linear law $\kappa_f(x) = Kx$,
with $K$ the LQR gain, is the one that connects back to dual-mode control: $\Xf$ is then a set
inside which the unconstrained feedback respects all the constraints, and the second mode of
\cref{fig:dualmode} is exactly the terminal law.

The choice involves a trade-off. The terminal constraint $x_{\ik{N}} \in \Xf$ is \emph{added}
to the problem; it is not required by the application, and like every constraint it can only
make the optimal cost worse. A larger terminal set therefore means a larger feasible set and a
lower cost, and one looks for the largest $\Xf$ that admits a law satisfying
\eqref{eq:terminalinv}. Under that condition the feasible sets also grow with the horizon:
$\mathcal{F}_N \subseteq \mathcal{F}_{N+1}$, because a sequence that reaches $\Xf$ in $N$ steps,
extended with one step of the terminal law, reaches it in $N+1$.

\begin{remark}[when feasibility fails]
The Building Energy Lab on the website makes the failure visible: lower the contractual power
limit while forcing several appliances into the same narrow window, and the optimizer reports
that no admissible plan exists. The constraints are contradictory; the algorithm is not at
fault. The remedy, when the application allows it, is to make the constraints that may fail
soft, as in \cref{sec:soft}.
\end{remark}

\begin{reading}
\refkey{Cannon\_2016}, section 2.4, ``Invariant Sets'' (pp.~21--24), and section 2.5,
``Controlled Invariant Sets and Recursive Feasibility'' (pp.~24--29): the proof above in the
same form, with the nested feasible sets. \refkey{Rawlings\_2017}, section 2.4.2,
``Stabilizing Conditions'' (pp.~114--120), for the assumptions on the terminal set in their
standard form, and section 2.6, ``Is a Terminal Constraint Set Necessary?'' (pp.~144--147), for
what happens without one.
\end{reading}
